\begin{figure}[htbp]
    \centering
    \begin{subfigure}[b]{0.48\textwidth}
        \centering
        % 導入純淨的正方形九宮格
        \includegraphics[width=\textwidth]{D:/Joker/Timoshenko/fig/2D_ssss/2d4s_vi_q4_ex/output_grid/17x17_ex_w1_9_grid.png}
        \caption{Navier 級數解析解 (1--9 階)}
    \end{subfigure}
    \hfill
    \begin{subfigure}[b]{0.48\textwidth}
        \centering
        % 導入純淨的正方形九宮格
        \includegraphics[width=\textwidth]{D:/Joker/Timoshenko/fig/2D_ssss/2d4s_vi_q4_ex/output_grid/17x17_FEM_w1_9_grid.png}
        \caption{Q4-SRI 有限元素數值解 (1--9 階)}
    \end{subfigure}
    
    \vspace{0.3cm} % 留下一點小空隙
    
    \begin{subfigure}[b]{0.48\textwidth}
        \centering
        % 導入純淨的正方形九宮格
        \includegraphics[width=\textwidth]{D:/Joker/Timoshenko/fig/2D_ssss/2d4s_vi_q4_ex/output_grid/17x17_ex_w10_15_grid.png}
        \caption{Navier 級數解析解 (10--15 階)}
    \end{subfigure}
    \hfill
    \begin{subfigure}[b]{0.48\textwidth}
        \centering
        % 導入純淨的正方形九宮格
        \includegraphics[width=\textwidth]{D:/Joker/Timoshenko/fig/2D_ssss/2d4s_vi_q4_ex/output_grid/17x17_FEM_w10_15_grid.png}
        \caption{Q4-SRI 有限元素數值解 (10--15 階)}
    \end{subfigure}

    % 在兩張大圖下方，居中放上一張共用的、扁平獨立的 Colorbar 橫圖
    \includegraphics[width=0.6\textwidth]{D:/Joker/Timoshenko/fig/2D_ssss/2d4s_vi_q4_ex/17x17_FEM_colorbar.png}
    
    \caption{正方形簡支 Mindlin 板前 15 階自由振動位移模態與下方共用色彩標尺（完美正方形矩陣排版）}
    \label{fig:modes_1_15_perfect}
\end{figure}